\documentclass[11pt]{article}
\usepackage{fontspec,amsmath,amssymb,amsthm,hyperref}
\usepackage[margin=1in]{geometry}
\setlength{\emergencystretch}{3em}
\title{Canonical Gaussian path law and implementation independence}
\author{LangevinLean mathematical exposition}
\date{}
\begin{document}\maketitle
\noindent Source module: \texttt{IPLA/AlgorithmLaw.lean}. Compile this proof document with XeLaTeX or LuaLaTeX. The accompanying HTML places the complete Lean declarations beside every proof.


The scalar Gaussian coordinate hypotheses uniquely determine the entire countable noise path. Independence of initialization then determines every output law, even across different underlying probability spaces.


\section{Construct the independent Gaussian noise path}


A noise coordinate is a pair \((k,a)\) consisting of iteration and state coordinate. Put the countable product law \(\bigotimes_{(k,a)}\mathcal N(0,1)\) on scalar paths \(w\). Packing sets \((\operatorname{pack}w)_k(a)=w(k,a)\); its measurability follows because every output coordinate is a measurable product projection, followed by the measurable Euclidean coordinate identification. The canonical vector path law is the pushforward of that scalar product law by packing.



The two probability instances are proved as well: each scalar Gaussian factor has mass one, hence the infinite product is a probability measure; pushing it through measurable packing preserves total mass. These are \texttt{canonicalFlatNoiseLaw\_isProbabilityMeasure} and \texttt{canonicalNoisePathLaw\_isProbabilityMeasure}, whose source appears in the complete-module panel.



\noindent Lean declarations: \texttt{NoiseCoordinate}, \texttt{canonicalFlatNoiseLaw}, \texttt{packNoisePath}, \texttt{packNoisePath\_measurable}, \texttt{canonicalNoisePathLaw}, \texttt{canonicalFlatNoiseLaw\_isProbabilityMeasure}, \texttt{canonicalNoisePathLaw\_isProbabilityMeasure}.

\section{The algorithm hypotheses identify the entire path law}


Mutual independence of the scalar family implies that its joint map has law equal to the infinite product of its individual coordinate laws. Since each law is \(\mathcal N(0,1)\),

\[\mathcal L((\xi_{k,a})_{k,a})=\bigotimes_{k,a}\mathcal N(0,1).\]

Composing this map with measurable packing and using the composition rule for pushforwards gives \(\mathcal L((\xi_k)_k)=\operatorname{canonicalNoisePathLaw}\). The argument proves independence of vector blocks from scalar independence and never assumes a grouping principle without justification.



\noindent Lean declarations: \texttt{GaussianInnovations.flatNoisePath\_law}, \texttt{GaussianInnovations.noisePath\_law}.

\section{Identify the common update and recursion}


The generic pointwise ULA update for \(V_N\) and step \(\gamma/N\) is \(z-(\gamma/N)\nabla V_N(z)+\sqrt{2(\gamma/N)}\xi\). Reassociating multiplication inside the square root makes it the scaled IPLA update. The generic noise-driven chain and \texttt{scaledChain} therefore agree by induction: their initial values match and each successor uses the same map on the same preceding state and noise vector.



\noindent Lean declarations: \texttt{pointwiseUpdate\_eq\_scaledUpdate}, \texttt{noiseDrivenChain\_eq\_scaledChain}.

\section{A common initial law gives common output laws}


Let \(F_k(z_0,\xi)\) be the measurable map that iterates the update \(k\) times. Initialization–path independence gives the joint input law \(\mu_0\otimes\mathcal G\), where \(\mathcal G\) is the canonical Gaussian path law. Hence

\[\rho_k=(F_k)_\#(\mu_0\otimes\mathcal G).\]

The pointwise/L\ensuremath{{}^{2}} identity identifies this measure with \texttt{laws}. The right-hand side depends only on \(U,\gamma,\mu_0\), so two Gaussian implementations with the same initial law give identical output laws. This identity needs no stationary diffusion or reference coupling.



\noindent Lean declarations: \texttt{laws\_eq\_canonicalChainLaw}, \texttt{laws\_eq\_of\_initialLaw\_eq}.

\end{document}
